\documentclass[12pt,a4paper]{article}
\usepackage[utf8]{inputenc}
\usepackage[T1]{fontenc}
\usepackage{geometry}
\usepackage{amsmath,amssymb,array,tabularx}
\geometry{margin=2.2cm}
\setlength{\parindent}{0pt}
\setlength{\parskip}{0.5em}
\title{oxygenation\_diagnostic\_mvp\\Detailed Diagnostic and Fitting-Integration Documentation}
\author{Synchronized to current module implementation}
\date{23 April 2026}

\begin{document}
\sloppy
\maketitle
\tableofcontents
\newpage

\section{Purpose}
The module \texttt{oxygenation\_diagnostic\_mvp.py} transforms blood-gas and tissue-sensor information into a probabilistic oxygenation assessment and alert category. It is designed to be explicit and auditable: each output can be traced to explicit equations and fixed parameters.

\section{Data model and entry point}
\subsection{Input dataclass}
Main input type:
\begin{center}
\texttt{OxygenationInput(po2, pco2, pH, temperature\_c, sensor\_po2, hemoglobin\_g\_dl=13.5, venous\_sat\_percent=75.0)}
\end{center}

Meaning of inputs:
\begin{itemize}
    \item \texttt{po2}: blood-gas PO$_2$,
    \item \texttt{pco2}: blood-gas PCO$_2$,
    \item \texttt{pH}: blood-gas acidity,
    \item \texttt{temperature\_c}: body temperature,
    \item \texttt{sensor\_po2}: local tissue/sensor PO$_2$,
    \item \texttt{hemoglobin\_g\_dl}: optional hemoglobin,
    \item \texttt{venous\_sat\_percent}: optional venous saturation.
\end{itemize}

\subsection{Main API}
\texttt{alert\_decision(oxygen\_input, yellow\_threshold=0.25, orange\_threshold=0.50, red\_threshold=0.75)}

Input guard: thresholds must satisfy
\[
0 \leq yellow \leq orange \leq red \leq 1.
\]
Otherwise a \texttt{ValueError} is raised.

\section{Mathematical functions used by diagnostics}
\subsection{Helper transforms}
\begin{itemize}
    \item \texttt{\_clamp(x)}: bounds values to a finite interval.
    \item \texttt{\_sigmoid(x)}: $\sigma(x)=\frac{1}{1+e^{-x}}$, with internal clipping of $x\in[-60,60]$ for numerical stability.
    \item \texttt{\_softmax(scores)}: $p_i=\frac{e^{s_i}}{\sum_j e^{s_j}}$, with max-shift stabilization.
\end{itemize}

\subsection{Bohr/CO$_2$/temperature correction}
\texttt{effective\_p50(pH, pco2, temperature\_c)} computes:
\[
P_{50,eff}=P_{50,std}\cdot10^{\beta_{pH}(pH-PH_{ref}) + \beta_{CO_2}\log_{10}(pCO_2/PCO2_{ref}) + \beta_T(T-T_{ref})}.
\]
Constants are imported from \texttt{krogh\_app/constants.py} and match the forward model constants.

\subsection{Hill saturation for estimated arterial saturation}
\texttt{hill\_saturation(po2, p50\_eff)}:
\[
S(P)=\frac{P^n}{P^n+P_{50,eff}^n},\quad n=2.7.
\]
In \texttt{alert\_decision}, this yields:
\[
\texttt{estimated\_sa\_percent}=100\cdot S(PO_2).
\]

\section{Risk-score construction}
\subsection{Feature-level risks}
Current feature transforms:
\begin{align*}
R_{po2} &= \sigma\left(\frac{55-po2}{12}\right), \\
R_{sensor} &= \sigma\left(\frac{35-sensor\_po2}{10}\right), \\
R_{acid} &= \sigma\left(\frac{7.32-pH}{0.06}\right), \\
R_{hypercapnia} &= \sigma\left(\frac{pco2-48}{8}\right), \\
R_{anemia} &= \sigma\left(\frac{11.5-Hb}{1.5}\right), \\
R_{venous} &= \sigma\left(\frac{68-Sv}{7}\right), \\
R_{gap} &= \sigma\left(\frac{(po2-sensor\_po2)-45}{15}\right).
\end{align*}

Temperature risk is the maximum of low- and high-temperature stress sigmoids:
\[
R_T = \max\left(\sigma\left(\frac{35-T}{0.8}\right),\;\sigma\left(\frac{T-38.5}{0.8}\right)\right).
\]

Compensation-gap term used later in state scoring:
\[
\texttt{compensation\_gap}=\texttt{clamp}\left(\frac{po2-sensor\_po2}{80}\right).
\]

\subsection{Weighted composite score}
Current weights in code:
\[
R = 0.22R_{po2}+0.24R_{sensor}+0.08R_{gap}+0.12R_{acid}+0.10R_{hypercapnia}+0.08R_{anemia}+0.06R_{venous}+0.02R_T.
\]
Then clamp to $[0,1]$.

Hard-floor overrides:
\begin{itemize}
    \item if \texttt{sensor\_po2 < 15} or \texttt{po2 < 25}, enforce $R\ge 0.90$,
    \item else if \texttt{sensor\_po2 < 25} or \texttt{po2 < 35}, enforce $R\ge 0.72$.
\end{itemize}

These overrides ensure severe hypoxic signals are not diluted by averaging in the weighted score.

\section{State probabilities and class assignment}
\subsection{Five-state framework}
States in current model:
\begin{enumerate}
    \item \texttt{normoxia},
    \item \texttt{intermediate\_oxygenation},
    \item \texttt{low\_oxygenation\_approaching\_critical},
    \item \texttt{hypoxia},
    \item \texttt{profound\_hypoxia}.
\end{enumerate}

\subsection{Score construction}
For each state, a Gaussian-like PO$_2$-band score is built around fixed centers and widths, then modified by:
\begin{itemize}
    \item direct sensor PO$_2$ heuristics,
    \item risk-score and compensation-gap contributions,
    \item additional boosts for very low sensor or blood-gas PO$_2$.
\end{itemize}

Finally \texttt{\_softmax} normalizes to probabilities.

\subsection{Confidence and certainty}
\begin{itemize}
    \item \texttt{confidence}: maximum state probability,
    \item \texttt{certainty}: gap between highest and second-highest probability (clamped to [0,1]).
\end{itemize}

\subsection{Alert level mapping}
Using user thresholds:
\begin{itemize}
    \item red if $R\ge red$,
    \item orange if $orange\le R < red$,
    \item yellow if $yellow\le R < orange$,
    \item green otherwise.
\end{itemize}

\section{Dominant-driver explanation logic}
\texttt{\_describe\_feature\_risks(...)} ranks feature risks and generates:
\begin{itemize}
    \item \texttt{dominant\_risk\_drivers}: top human-readable drivers,
    \item \texttt{driver\_summary}: sentence-style summary for GUI/report text.
\end{itemize}

This is central for interpretability because the user sees not only alert color, but also the main contributors.

\section{Return payload structure}
\texttt{alert\_decision(...)} returns a dictionary including:
\begin{itemize}
    \item \texttt{predicted\_state}, \texttt{risk\_score}, \texttt{alert\_level},
    \item \texttt{confidence}, \texttt{certainty},
    \item five primary probabilities,
    \item backward-compatible alias probabilities for existing report/GUI paths,
    \item \texttt{estimated\_sa\_percent}, \texttt{p50\_eff},
    \item \texttt{feature\_risks}, \texttt{dominant\_risk\_drivers}, \texttt{driver\_summary},
    \item normalized \texttt{input} echo.
\end{itemize}

\section{How this module is used in the full application}
\subsection{Call path}
\begin{enumerate}
    \item User enters diagnostic inputs in GUI diagnostic tab.
    \item GUI creates \texttt{OxygenationInput} and calls diagnostic engine wrapper.
    \item Wrapper calls \texttt{alert\_decision(...)}.
    \item GUI displays probabilistic state, risk score, certainty, alert, and driver summary.
\end{enumerate}

\subsection{Coupling to Krogh reconstruction}
After diagnostic evaluation, the GUI can run inverse reconstruction:
\begin{itemize}
    \item venous saturation is transformed to a venous PO$_2$ target,
    \item reconstructor fits perfusion, mitoP50, and relative metabolic demand,
    \item hidden hypoxic burden and radius-conditioned variants (30/50/100 \textmu m) are reported,
    \item uncertainty and identifiability summaries are attached.
\end{itemize}

Therefore, diagnostics and reconstruction should be interpreted jointly.

\section{Function-to-code mapping table}
\begin{center}
\renewcommand{\arraystretch}{1.1}
\small
\begin{tabularx}{0.98\textwidth}{|>{\raggedright\arraybackslash}p{0.28\textwidth}|>{\raggedright\arraybackslash}p{0.32\textwidth}|>{\raggedright\arraybackslash}X|}
\hline
Function & Inputs & Output role \\
\hline
\texttt{effective\_p50} & pH, pCO2, temperature & Shifts oxygen affinity for current blood-gas context \\
\hline
\texttt{hill\_saturation} & po2, optional p50\_eff & Estimated saturation, used for \texttt{estimated\_sa\_percent} \\
\hline
\texttt{\_state\_probabilities} & risk\_score, compensation\_gap, sensor\_po2, po2 & Five-state probability map \\
\hline
\texttt{\_describe\_feature\_risks} & feature risk dictionary & Dominant-driver labels and narrative summary \\
\hline
\texttt{alert\_decision} & \texttt{OxygenationInput}, thresholds & Main diagnostic payload for GUI/report layers \\
\hline
\end{tabularx}
\normalsize
\end{center}

\section{Output and report integration}
\subsection{Where diagnostic outputs are stored}
The diagnostic result dictionary is cached in GUI runtime as \texttt{last\_diagnostic\_result}. If reconstruction runs, \texttt{last\_krogh\_reconstruction} is also cached.

\subsection{Export}
Via \texttt{krogh\_app/persistence.py}:
\begin{itemize}
    \item diagnostic report export as JSON or CSV,
    \item publication report export as TXT/TEX/PDF,
    \item automatic inclusion of dominant drivers, uncertainty summary, hidden burden summary, and radius sensitivity summary.
\end{itemize}

\section{Interpretation guidance for users}
\begin{itemize}
    \item The diagnostic score is a transparent heuristic model, not a trained clinical predictor.
    \item High confidence with low certainty means one state dominates only weakly over the second-best state.
    \item Strong alert levels should be interpreted with reconstruction assumptions (fit warnings, boundary hits, radius sensitivity).
    \item Repeated measurements and trend interpretation are preferred over isolated snapshots.
\end{itemize}

\section{Conclusion}
\texttt{oxygenation\_diagnostic\_mvp.py} provides a mathematically explicit probabilistic oxygenation interpretation layer that is tightly integrated with the mechanistic Krogh reconstruction workflow. The implementation is intentionally transparent and suitable for scientific audit and reproducible reporting.

\end{document}
